
\documentclass[11pt, a4paper]{article}
\usepackage{graphicx}
\usepackage{listings}
\usepackage{xcolor}
\usepackage{hyperref}
\usepackage{amsfonts}
\usepackage[margin=0.5in,top=1in,bottom=1in]{geometry}

\title{The $\mathbb{C}$ Index Property}
\author{Amlal El Mahrouss\\amlal@nekernel.org}
\date{February 2026}

\begin{document}
\bf
\maketitle
\begin{center}
	\rule[0.01cm]{17cm}{0.01cm}
\end{center}
\abstract{
	This paper defines the property of $\mathbb{C}$ Indices such that $n \in \mathbb{R}$ for all $\Omega_{n} \in \mathbb{C}$. Definitions and properties are stated in this paper.
}
\begin{center}
	\rule[0.01cm]{17cm}{0.01cm}
\end{center}

\section{Definitions of $\mathbb{C}$ Indices}

\subsection{Definition of $\Omega$}

Let $\Omega_{n}$ be a set defined as $\Omega \in \mathbb{C}$ and the following indices $\beta \in \mathbb{C}$:

\begin{equation}
	\Omega_{n} = (\beta_{0} \times \beta_{n})
\end{equation}

\subsection{Definition of $\alpha$}

Let $\alpha$ denote $\Omega_{n}$ with a remainder of $\mathbb{X}$ defined as the result of $\beta_{n}$.
\begin{equation}
	\alpha = \Omega_{n} \cup \mathbb{X}
\end{equation}

\section{Properties of $\mathbb{C}$ Indices}

\subsection{Properties of $\Omega$}

\begin{enumerate}
	\item $\Omega_{n} \cup \mathbb{X}$ such that $\mathbb{X} = \beta_{n}$.
\end{enumerate}

\subsection{Properties of $\alpha$}

\begin{enumerate}

	\item $\mathbb{X} = \beta_{n}$ such that $\beta_{n} \in \mathbb{C}$
	
\end{enumerate}

\end{document}
